\section{Conclusion and Future Work}
\label{sec:conclusion}

We have presented an initial implementation of relational e-matching inside an
SMT solver, identifying \emph{backtracking} as the central obstacle. Our
synthetic benchmarks demonstrate that the relational approach can substantially
outperform the conventional top-down matcher. Our results on Verus benchmarks
suggest that it is not that simple to translate this asymptotic advantage into a
practical win on real workloads. Closing the remaining gap will require a
combination of better data structures, better query planning, and a tighter
integration with the rest of the solver. We sketch the most promising directions
below.

\paragraph{Better backtrackable index maps.}
Our current implementation maintains
indices as persistent maps, which is correct but leaves substantial performance
on the table. Other datastructures are backtracked using trail-based undo logs, or
amortized rebuilding schemes. Bringing these to bear on the
relational indices would directly translate into reduced per-decision overhead.

\paragraph{Understanding what real benchmarks need.}
Synthetic benchmarks let us showcase the asymptotic advantage of relational
e-matching. We do not yet
have a precise characterization of which patterns and which workloads benefit,
and which do not. A more careful study --- profiling pattern shapes, e-graph sizes,
and the distribution of trigger applications across decision levels --- is needed.

\paragraph{Multi-query optimization.}
Today our implementation compiles each trigger into its own query plan and
runs each plan independently. SMT solvers, however, evaluate hundreds or
thousands of triggers against the same e-graph; many of these triggers share
sub-patterns. The database community has long studied multi-query optimization,
and Z3's code trees~\cite{ematchingz3} can be viewed as exactly this kind of
sharing applied to top-down e-matching. We expect that a relational analogue
--- sharing common sub-joins across triggers, or jointly planning batches of
triggers --- would yield substantial savings on real workloads.

\paragraph{Heuristics inside the query planner.}
Relational e-matching opens the door to a much richer space of optimizations
than top-down matching, because the planner now has visibility into what work
will be done. Many SMT-specific heuristics fit naturally into this framework.
For instance, \emph{relevancy filtering}~\cite{relevancy} --- restricting
matches to terms the SAT engine currently considers active --- can be
expressed as an additional selection in the join plan, and a planner that
knows about relevancy can prune entire join paths early. Other heuristics, such as preferring patterns
that have produced useful instantiations in the past, or biasing the join
order toward terms involved in recent conflicts, similarly fit as cost-model
inputs to the planner.

% \paragraph{Outlook.}
Indexing techniques have proven helpful in SMT
performance: code trees and inverted path indices~\cite{ematchingz3} are widely
credited with making large-scale e-matching tractable in Z3. Relational
e-matching is, in a precise sense, the most general indexing technique
available --- it subsumes pattern compilation as a special case and grants the
planner full freedom to reorder and share work. With backtrackability solved
and the right query-planning heuristics in place, we believe relational
e-matching has a real path to becoming a practical and broadly useful
component of modern SMT solvers.
